\section*{Chapter \ref{ch:g7:project-quiz-game} --- Project: A Quiz Game}

\begin{solution}{exo:g7:project-quiz-game:1}
\code{"paris"} and \code{"new york"}: the space inside stays, only the spaces
at the ends go.
\end{solution}

\begin{solution}{exo:g7:project-quiz-game:2}
\code{eight} is not \code{8}, and \code{hot} is wrong; the other eight answers
are right once normalised, \code{" Honey "} included: \code{Your score: 8 out
of 10}.
\end{solution}

\begin{solution}{exo:g7:project-quiz-game:3}
The typed answer is normalised once, and tested for \code{skip} first:

\omcode{code/grade-7/06-project-quiz-game/exercises/skip.py}{5}{18}{Skipping a
question.}\label{lst:g7:project-quiz-game:sol-skip}
\end{solution}

\begin{solution}{exo:g7:project-quiz-game:4}
An \omterm{def:g7:lists:list}{item} of \code{answers} was deleted after the fourth: from question 5 on,
each question is compared with the answer to the next one, and the tenth
question has no answer at all, so \code{answers[9]} is past the end. The
check \code{if len(questions) != len(answers):} of version 1 would have
printed \code{Error: 10 questions but 9 answers}.
\end{solution}

\begin{solution}{exo:g7:project-quiz-game:5}
In \code{main}, after the result line:
\texttt{print(round(score / len(questions) * 100), "\%")}, which prints
\texttt{80 \%} for the game of \cref{out:g7:project-quiz-game:v5}.
\end{solution}

\begin{solution}{exo:g7:project-quiz-game:6}
The second question is asked even after a right first answer, and its half
point is added to the 1. The \omterm{def:g7:functions:function}{function} must return as soon as the first
answer is right, as \code{ask} of version 4 does with its first
\code{return 1}.
\end{solution}

\begin{solution}{exo:g7:project-quiz-game:7}
\code{questions.py} gets two more \omterm{def:g7:lists:list}{lists}, \code{MEDIUM_QUESTIONS} and
\code{MEDIUM_ANSWERS}; \code{choose_level} accepts \code{"medium"} in its
\code{while} \omterm{def:g4:algorithms-repeat-and-decide:condition}{condition}; \code{main} gets one more \code{elif} to pick the
medium \omterm{def:g7:lists:list}{lists}. \code{normalise}, \code{ask} and \code{play} do not change: they
work on any pair of \omterm{def:g7:lists:list}{lists}.
\end{solution}

\begin{solution}{exo:g7:project-quiz-game:8}
One line for the level, at most three per question (the question, the second
try, the answer shown), and the result line: $1 + 10 \times 3 + 1 = 32$. A
test of the chapter plays a game with every answer wrong and counts 32. A
player who keeps typing a wrong level is asked again and again, one more line
each time, without end.
\end{solution}

\begin{solution}{exo:g7:project-quiz-game:9}
\omcode{code/grade-7/06-project-quiz-game/exercises/quit.py}{5}{17}{Stopping
early.}\label{lst:g7:project-quiz-game:sol-quit}

The \omterm{def:g6:variables-and-input:value}{variable} \code{playing} becomes \code{False} at \code{quit}, which ends
the \omterm{def:g4:algorithms-repeat-and-decide:loop}{loop} before the next question.
\end{solution}

\begin{solution}{exo:g7:project-quiz-game:10}
A running count of right first answers, set back to 0 by any other result,
and the best count so far:

\omcode{code/grade-7/06-project-quiz-game/exercises/streak.py}{1}{15}{The
longest streak.}\label{lst:g7:project-quiz-game:sol-streak}

It prints 3, for questions 3 to 5.
\end{solution}

\begin{solution}{exo:g7:project-quiz-game:11}
The highest is 10, every answer right at once; the lowest is 0. With $a$ right
first answers and $b$ right second answers, the score is $a + b/2$ with
$a + b \le 10$. Any number of halves $k$ from 0 to 20 can be made: if $k$ is
even, take $a = k/2$ and $b = 0$; if $k$ is odd, take $a = (k-1)/2$ and
$b = 1$, and then $a + b \le 10$ since $k \le 19$. So the 21 scores 0, 0.5,
1, \dots, 10 can all be reached. A \omterm{def:g4:programs-and-operating-system:program}{program} that tries every $a$ and $b$ (the
\omterm{def:g3:files-and-folders:file}{file} \code{exercises/scores.py}) finds the same 21.
\end{solution}

\begin{solution}{exo:g7:project-quiz-game:12}
Each answer becomes a \omterm{def:g7:lists:list}{list}, and \code{in} tests whether the typed answer is one
of them:

\omcode{code/grade-7/06-project-quiz-game/exercises/many_answers.py}{5}{18}{Several
accepted answers.}\label{lst:g7:project-quiz-game:sol-many}
\end{solution}
